\subsection{Test set primary results}
The frozen, validation-selected checkpoints are evaluated on 4,107 images, 240 patients, and four centers. The order is VCRE-Fib, w/o Weak Loc., w/o View, and SFibAI (Table~\ref{tab:comparison}), with lower $\risk$ indicating better performance. Main Tables~\mainref{tab:main}--\mainref{tab:auxiliary} report the primary, image, and spatial outcomes; Appendix~\ref{sec:appendix-patient} gives the detailed patient-level metrics.
\inputtable{tab04_comparison}

\paragraph{Metric definitions.}
Let $e=|\hat y-y|$ and let $d$ be the absolute difference in four-level grades with thresholds 0.5, 1.5, and 2.5. The locked per-item risk is
\[
\operatorname{COR}=0.35e/3.5+0.15\mathbf{1}(e>0.3)+0.15\mathbf{1}(e>0.5)+0.25d/3+0.10\mathbf{1}(d\ge2).
\]
Patient-max and patient-median aggregate reference and predicted scores separately. Center-balanced COR uses square-root patient-count weights. TMAE is $\operatorname{mean}\max(e-0.5,0)$ and severe-error rate is $\operatorname{mean}\mathbf{1}(e>1.0)$, distinct from the two-grade error term inside COR.
\subsection{Detailed patient-level grading}\label{sec:appendix-patient}
Under patient-max and patient-median aggregation, the full model retains the lowest COR (0.200528 and 0.154344), alongside the lowest median MAE (0.339273; Table~\ref{tab:patient}) and center-balanced COR (0.177453; main-text Table~\mainref{tab:main}). Some thresholded accuracies still favor variants: w/o Weak Loc.\ has higher patient-max $\pm0.5$ accuracy (67.50\%), and w/o View has higher median four-grade accuracy (72.08\%). The full model's advantage thus spans overall risk and multiple grading measures.
\inputtable{tab06_patient}

\subsection{Additional grading metrics and center-level results}
Per-center estimates are descriptive; one center has two patients. Test/val exports retain per-grade support/recall and all COR terms. Four-grade probabilities sum the 36 bins; image AUPRC uses average precision and ECE uses ten bins. Class-wise AUROC/AUPRC are reported as N/A when either positive or negative examples are absent. Patient probabilities are not aggregated.
\inputtable{tab14_additional}
\inputtable{tab15_centers}
\FloatBarrier
